\emph{Consolidated reference, v4. This document has been revised across three rounds of external review: an evidentiary-accuracy pass, a systems-safety pass (which added Parts Eleven through Fourteen), and a purpose-and-standard-of-success pass (which added Part Two). A dedicated section at the end (Part Fifteen) lists points raised across all three reviews that remain open to debate rather than settled corrections --- treat that section as part of the working discussion, not as resolved.}

\textbf{A note on how to read this document.} Not every sentence here carries the same evidentiary weight, and v1's biggest flaw was not distinguishing between them. Going forward, treat statements in five rough categories: \textbf{definitions} (a terminological choice this document adopts for clarity --- not true or false, just useful or not); \textbf{framework propositions} (an analytical lens this document proposes --- useful if it clarifies thinking, not true by citation); \textbf{mechanistic claims} (how a system actually works, e.g. next-token training objectives --- these are about as solid as anything here gets); \textbf{empirical claims} (a measured finding --- these date fast, vary by benchmark and methodology, and should be treated as illustrative rather than current fact unless independently re-verified at the time of use); and \textbf{normative or legal claims} (a recommended posture or a jurisdiction-specific proposition --- never a substitute for professional legal advice). Where a claim in v1 blurred one of these into another --- most often, presenting a framework proposition or a single illustrative example as if it were an established empirical law --- that's exactly what this revision tries to fix.

\textbf{A note on this version (v4).} A third review argued that the document, as it stood, could read primarily as an argument for finding reasons \emph{not} to deploy AI systems --- a long, careful catalogue of what can go wrong with no stated positive objective against which "good enough" could ever be declared. That's a fair critique of an omission rather than a correction of an error, and it's addressed by a new \textbf{Part Two: Purpose and Standard of Success}, placed immediately after Part One so that everything downstream is read against a stated goal rather than against an implicit standard of perfection. Adding it shifts the numbering: what was v3's Part Two ("Define the Job") is now Part Three, and each part after it through v3's Part Nine ("Threads That Tie Everything Together") shifts up by one, to Part Four through Part Ten. Parts Eleven through Fifteen --- the error-propagation material added in v3 --- keep their original numbers unchanged, since v3's numbering already had a gap at Part Ten that this shift happens to close exactly.

\textbf{A note on this version specifically (v3).} A second review, from a systems-safety rather than an evidentiary-accuracy angle, identified a genuine gap rather than a factual error: the document had substantial coverage of where errors originate and why verification is hard, but no explicit theory of how an erroneous, uncertain, or incomplete state actually \emph{travels} through a socio-technical system and becomes --- or fails to become --- real-world harm. Parts Eleven through Fourteen add that theory, organized around three planes that recur throughout them: the \textbf{epistemic} plane (what is known, unknown, wrong, disputed, or omitted), the \textbf{causal} plane (how that state propagates, amplifies, disappears, or gets embedded), and the \textbf{operational} plane (what prevents harm, detects failure, enables recovery, and assigns responsibility). The two views are complementary: the tool- and output-centric parts are largely about a single output or a single tool evaluated in isolation, and Parts Eleven through Fourteen are about what happens once that output leaves the tool and enters an organization's actual working life.
